\section{问题一：单日购电与储能协同优化}
\subsection{决策变量与目标函数}
在日初已知全天输入的条件下，以144个十分钟区间为统一决策时域。对第$k$个区间，设置购电量$g_k$、充电量$c_k$、放电量$d_k$和光伏消纳量$u_k$，并以$E_k$表示区间开始时的储电量。购电费用最小化目标为
\begin{equation}
  \min J=\sum_{k=0}^{143}p_kg_k.
  \label{eq:objective}
\end{equation}
储能会增加转换损耗，却可能替代高价时段的购电，故本文直接优化电价加权费用。

\subsection{约束条件}
\subsubsection{供需平衡}
各区间购电、储能放电与光伏消纳共同供给负载和充电需求，满足
\begin{equation}
  g_k+d_k+u_k=\ell_k+c_k,\qquad k=0,\ldots,143.
  \label{eq:balance}
\end{equation}
购电和光伏消纳的边界为
\begin{equation}
  g_k\ge0,\qquad 0\le u_k\le v_k.
  \label{eq:pv}
\end{equation}
其中$v_k-u_k$为弃光量。当光伏富余而储能容量或功率不足时，允许弃光可维持供需平衡；本模型不设售电变量。

\subsubsection{储能递推与初末状态}
母线侧充入$c_k$时，电池内部增加$\eta_{\mathrm c}c_k$；母线侧放出$d_k$时，内部减少$d_k/\eta_{\mathrm d}$。由此有
\begin{equation}
  E_{k+1}=E_k+\eta_{\mathrm c}c_k-\frac{d_k}{\eta_{\mathrm d}}
  =E_k+0.9c_k-\frac{d_k}{0.9}.
  \label{eq:state}
\end{equation}
全部边界状态均须满足运行范围，并使日末恢复至日初水平：
\begin{equation}
  1200\le E_k\le10800,\quad k=0,\ldots,144;
  \qquad E_0=E_{144}=6000.
  \label{eq:bounds}
\end{equation}
初末状态相等，使费用不因净消耗日初已有电量而被低估。

\subsubsection{功率限制与充放电互斥}
最大母线侧功率5\,000 kW对应的十分钟电量上限为
\begin{equation}
  M=P_{\max}\Delta t=\frac{5000}{6}\ \mathrm{kWh}.
  \label{eq:M}
\end{equation}
引入二进制变量$z_k$，将充放电互斥写为
\begin{equation}
  0\le c_k\le Mz_k,\qquad
  0\le d_k\le M(1-z_k),\qquad z_k\in\{0,1\}.
  \label{eq:mutex}
\end{equation}
$z_k=1$允许充电，$z_k=0$允许放电，待机时两者均为零。此处的$M$由设备参数确定，避免使用过大的任意常数。

\subsection{模型求解}
式\eqref{eq:objective}--\eqref{eq:mutex}构成混合整数线性规划。变量按$(g,c,d,u,E,z)$排列，共有721个连续变量和144个二进制变量；供需平衡与状态递推形成288个等式，互斥条件形成288个不等式，其余条件由变量边界给出。

采用SciPy 1.17.1的\texttt{milp}接口调用HiGHS求解\cite{scipyMilp,highs}，步骤见图\ref{fig:workflow}。求解后使用未舍入序列重新计算供需平衡、储能变化和费用，再汇总指定时段结果。本次目标值和求解器下界均为35\,126.948589元，相对最优间隙为0，在求解器数值精度下达到本模型的报告最优值。
\begin{figure}[H]
  \centering
  \PaperGraphic[width=0.98\linewidth]{figures/q1_workflow.pdf}
  \caption{问题一的模型建立与求解流程}
  \label{fig:workflow}
\end{figure}

\subsection{计算结果}
全天购电量为59\,482.699 kWh，购电费用为35\,126.949元。指定区间购电量与四小时充放电量分别见表\ref{tab:q1-purchase}、表\ref{tab:q1-battery}。表中仅在展示时保留三位小数，全部检验使用未舍入值。
% Source: q1-a-formal-003/domain_result.json; SHA-256: 52dd1f9efb6136e00e56114a9097e3f330c2c7afa4724af24b3e7fd270f1b0d4
\begin{table}[H]
\centering
\small
\caption{问题一指定区间购电量及全天购电结果（电量单位：kWh）}
\label{tab:q1-purchase}
\begin{tabular}{lrlrlr}
\toprule
时间段 & 购电量 & 时间段 & 购电量 & 时间段 & 购电量 \\
\midrule
10:00--10:10 & 0.000 & 12:00--12:10 & 480.412 & 14:00--14:10 & 0.000 \\
16:00--16:10 & 445.432 & 18:00--18:10 & 531.894 & 20:00--20:10 & 0.000 \\
\midrule
全天购电量 & \multicolumn{2}{l}{59\,482.699} & 全天购电费（元） & \multicolumn{2}{r}{35\,126.949} \\
\bottomrule
\end{tabular}
\end{table}
\begin{table}[H]
\centering
\small
\caption{问题一四小时充放电量与初末储电量（单位：kWh）}
\label{tab:q1-battery}
\begin{tabular}{lrrlrr}
\toprule
时间段 & 充电量 & 放电量 & 时间段 & 充电量 & 放电量 \\
\midrule
0:00--4:00 & 4\,500.000 & 0.000 & 4:00--8:00 & 833.333 & 6\,365.841 \\
8:00--12:00 & 4\,787.964 & 1\,702.997 & 12:00--16:00 & 5\,286.035 & 91.101 \\
16:00--20:00 & 0.000 & 5\,780.132 & 20:00--24:00 & 5\,333.333 & 2\,859.868 \\
\midrule
\multicolumn{3}{l}{0:00储电量：6\,000.000} & \multicolumn{3}{r}{24:00储电量：6\,000.000} \\
\bottomrule
\end{tabular}
\end{table}


\begin{figure}[H]
  \centering
  \PaperGraphic[width=0.96\linewidth]{figures/q1_dispatch.pdf}
  \caption{问题一的十分钟购电与充放电计划}
  \label{fig:dispatch}
\end{figure}
图\ref{fig:dispatch}中充电绘为正值、放电绘为负值，各值表示对应十分钟区间的电量。四小时汇总中充、放电量可同时为正，原因是设备在不同区间切换运行方式。例如4:00--8:00累计充电833.333 kWh、放电6\,365.841 kWh，而逐区间仍满足互斥条件。

\subsection{调度特点与能量收支}
结合图\ref{fig:input}、图\ref{fig:dispatch}及储电量轨迹图\ref{fig:soc}，可将调度过程理解为以下几个阶段。

0:00--4:00电价较低，储能累计充电4\,500.000 kWh，内部储电量由6\,000增加至10\,050 kWh，为随后时段供电。日间光伏逐步超过负载，8:00--16:00以净充电为主，16:00储电量达到10\,800 kWh上界。16:00--20:00光伏衰减且电价处于较高水平，储能累计放电5\,780.132 kWh，20:00储电量降至4\,377.631 kWh。此后的购电与充放电同时受到电价变化和日末恢复条件影响，最终回到6\,000 kWh。
\begin{figure}[H]
  \centering
  \PaperGraphic[width=0.96\linewidth]{figures/q1_soc.pdf}
  \caption{问题一的储电量轨迹与运行边界}
  \label{fig:soc}
\end{figure}

若低价时段购入1 kWh用于充电，完整充放电后最多向母线释放0.81 kWh。仅从购电价差看，当低价$p_{\mathrm{low}}$与高价$p_{\mathrm{high}}$满足$p_{\mathrm{low}}<0.81p_{\mathrm{high}}$时，跨时段转移才可能覆盖转换损耗；实际转移量还受到光伏、负载和设备边界限制。这解释了模型为何根据全天条件联合安排动作。

全天充电量为20\,740.666 kWh，放电量为16\,799.940 kWh。对式\eqref{eq:state}求和，并代入$E_{144}=E_0$，得到
\begin{equation}
  \sum_k d_k=\eta_{\mathrm c}\eta_{\mathrm d}\sum_k c_k
  =0.81\sum_k c_k.
  \label{eq:roundtrip}
\end{equation}
因此，充放电转换损耗为3\,940.727 kWh。再对供需平衡式求和，有
\begin{equation}
  \sum_k g_k=\sum_k\ell_k-\sum_k u_k+\sum_k(c_k-d_k).
  \label{eq:dailybalance}
\end{equation}
本次弃光量为零，55\,482.836 kWh光伏预测电量全部消纳。负载与光伏之差约55\,541.972 kWh，加上储能损耗即为全天购电量。表\ref{tab:q1-purchase}中三个指定区间购电量为零，仅表示对应区间由光伏与储能覆盖需求。
\FloatBarrier
